\documentclass[10pt,a4paper]{article}
\usepackage[margin=2.2cm]{geometry}
\usepackage{fontspec}
\setmonofont{DejaVu Sans Mono}[Scale=0.84]
\newfontfamily\unifont{DejaVu Serif}[Scale=0.9]
\usepackage{amsmath,amssymb,mathtools}
\setcounter{MaxMatrixCols}{20}
\usepackage{longtable,booktabs,array}
\usepackage{fancyvrb}
\usepackage[dvipsnames]{xcolor}
\usepackage[colorlinks=true,linkcolor=NavyBlue,urlcolor=NavyBlue]{hyperref}
\usepackage{microtype}
\usepackage{needspace}
\usepackage{tocloft}\setlength{\cftbeforesecskip}{0.45em}
\setlength{\parindent}{0pt}\setlength{\parskip}{0.55em}
\setlength{\emergencystretch}{8em}
\renewcommand{\arraystretch}{1.15}
\allowdisplaybreaks
\begin{document}
\title{The coupled field equations of the 7+1 model and their solutions in the three gravity regimes: symbols and their definitions\\[0.4em]\large A dictionary of every symbol used for the action term by term, the coupled field equations, the radion, and the accurate approximate solutions in the primordial, modern and far-future regimes — sections 20–23 of the notebook 7+1.nb}
\author{Repository: github.com/once-ere/7plus1 — prepared for Nash\\ Notebook: 7+1.nb (Wolfram Language 15.0.1)\\ Document generated by tools/make\_docs.py on 2026-10-27}
\date{}
\maketitle
\begin{abstract}
For each symbol this dictionary gives the mathematical notation, the name used in the Mathematica notebook 7+1.nb (where there is one), the defining equation and its meaning, grouped by topic. Longer defining equations follow each table as numbered displayed formulas. Every equation needed to define a symbol is written in this document.
\end{abstract}
\tableofcontents
\Needspace{12\baselineskip}\section*{1. Conventions}\addcontentsline{toc}{section}{1. Conventions}
\begin{itemize}
\item Units: $\hbar=c=1$. Sections 21 and 22.1 of the notebook (regime I and the capture of the auxiliary space) use the Gauss–Bonnet length $\ell_2=\sqrt{\alpha_2}=1$ as unit of length and time; sections 22.2–23 use GeV, s, K, Gyr where stated.
\item Coordinates $X=(x_0,x_1,\dots,x_7)$: $x_0$ is time (the direction of the octonion unit $e_0=1$), $x_1,x_2,x_3$ ordinary space, $x_4,\dots,x_7$ the four auxiliary space-like coordinates; on the world brane $x=(x_0,x_1,x_2,x_3)$, and $y=(x_4,\dots,x_7)$.
\item Indices: $\mu,\nu,\rho,\sigma,\lambda,\alpha,\beta,\gamma,\delta=0,\dots,7$ coordinate indices (on the brane $0,\dots,3$); $A,B,C=0,\dots,7$ frame indices; $k,l=4,\dots,7$ auxiliary directions; $i,j=1,2,3$ ordinary directions; repeated indices are summed.
\item A prime is $d/dx_0$; $\tau=x_0-x_s$ is the time since the singularity $x_s$; "ln" is the natural logarithm.
\item Signature: the bulk has one time-like and seven space-like dimensions; in xAct the list $\{7,1,0\}$ declares exactly this (seven space-like, one time-like, no degenerate dimension). The frame metric $\eta=\mathrm{diag}(+1,-1,\dots,-1)$ (\texttt{η17}) is simply the squares of the octonion basis, $\eta_{AB}=\mathrm{Re}(e_Ae_B)$.
\end{itemize}
\textbf{The signature: $\{7,1,0\}$ and $\mathrm{diag}(+1,-1,\dots,-1)$.} The 7+1 spacetime has one time-like dimension ($x_0$, the direction of the octonion unit $e_0=1$) and seven space-like dimensions ($x_1,x_2,x_3$ ordinary space, $x_4,\dots,x_7$ auxiliary space). Two different objects describe this, and neither may be replaced by the other. (i) The tensor package xAct, used in section 12 of 7+1.nb exactly as in the author’s notebook Einstein-Lovelock-7+1-Nash.nb, declares the metric with \texttt{DefMetric[\{7, 1, 0\}, g71[-a, -b], CD, …]}: the list $\{7,1,0\}$ declares seven space-like dimensions, one time-like dimension and no degenerate dimension — the numbers of entries $+1$ and $-1$ of $\eta_{71}=\mathrm{diag}(-1,+1,\dots,+1)$ (\texttt{η71}), the same spacetime written in the sign convention of 4+4.nb (space-like $=+1$, time-like $=-1$). (ii) The matrix $\eta=-\eta_{71}=\mathrm{diag}(+1,-1,\dots,-1)$ (\texttt{η17}) is simply the squares of the octonion basis: $e_0^2=+1$ for the unit $e_0=1$ (time) and $e_i^2=-1$ for the seven imaginary units $e_1,\dots,e_7$ (space); more generally it is the real part of the product of two basis units, and it gives the quadratic form of the octonions:
\begin{gather*}
\eta_{AB}=\mathrm{Re}(e_Ae_B),\qquad e_0^2=+1,\qquad e_i^2=-1\ \ (i=1,\dots,7),\\ \mathrm{Re}(x\,x)=x_0^2-x_1^2-\cdots-x_7^2=\eta_{AB}\,x^Ax^B\qquad (x=x^Ae_A\in\mathbb O).
\end{gather*}
$\eta$ gives the signs of the Clifford relation $\Gamma_A\Gamma_B+\Gamma_B\Gamma_A=2\,\eta_{AB}\,1_{16}$ of the real 16×16 Dirac matrices: time squares to $+1_{16}$ and each of the seven space directions to $-1_{16}$, exactly like the octonion units (in the notation $Cl(p,q)$, with $p$ generators squaring to $+1$ and $q$ to $-1$, the algebra $Cl(1,7)\cong\mathbb R(16)$). It also gives the signs of the metric components: $g_{00}>0$ along the time-like direction, negative squares along the space-like directions. Every component computed by 7+1.nb follows from the component matrix of the metric; the declared list fixes for xAct the sign of the determinant of the metric, \texttt{SignDetOfMetric[g71]} $=-1$, which is the same in both sign conventions, since $\det \eta=\det \eta_{71}=-1$. 7+1.nb checks \texttt{SignatureOfMetric[g71]} $=\{7,1,0\}$, \texttt{SignDetOfMetric[g71]} $=-1$, $\mathrm{Re}(e_Ae_B)=\eta_{AB}$ and $\mathrm{Re}(x\,x)=\eta_{AB}x^Ax^B$.

\Needspace{12\baselineskip}\section*{2. Spacetime, curvature and the Lovelock tensors}\addcontentsline{toc}{section}{2. Spacetime, curvature and the Lovelock tensors}
{\small
\begin{longtable}{>{\raggedright\arraybackslash}p{\dimexpr 0.1700\linewidth-2\tabcolsep\relax}>{\raggedright\arraybackslash}p{\dimexpr 0.2100\linewidth-2\tabcolsep\relax}>{\raggedright\arraybackslash}p{\dimexpr 0.3400\linewidth-2\tabcolsep\relax}>{\raggedright\arraybackslash}p{\dimexpr 0.2800\linewidth-2\tabcolsep\relax}}
\toprule
\textbf{symbol} & \textbf{Mathematica} & \textbf{definition} & \textbf{meaning} \\
\midrule\endfirsthead
\toprule
\textbf{symbol} & \textbf{Mathematica} & \textbf{definition} & \textbf{meaning} \\
\midrule\endhead
\endfoot
\bottomrule\endlastfoot
$g_{\mu\nu}$, $g^{\mu\nu}$, $|g|$ & \texttt{curvature[g]["g"]}, \texttt{["ginv"]}; \texttt{gCos} (cosmological) & $ds^2=g_{\mu\nu}dX^\mu dX^\nu$ & metric, inverse, $|\det g|$ \\
$\{{}^{\lambda}_{\mu\nu}\}$ & \texttt{curvature[g]["chr"]} & (2.1) & Christoffel symbols \\
$R^\rho{}_{\sigma\mu\nu}$, $R_{\sigma\nu}$, $R$ & \texttt{["riem"]}, \texttt{["ric"]}, \texttt{["R"]} of \texttt{curvature[g]} & (2.1) & Riemann tensor, Ricci tensor, scalar curvature \\
$G^\mu{}_\nu$ & \texttt{curvature[g]["G"]} & $R^\mu{}_\nu-\tfrac12\delta^\mu_\nu R$ & Einstein tensor \\
$\delta^{\mu_1\cdots\mu_n}_{\nu_1\cdots\nu_n}$ & \texttt{kδ} (xAct), pair sets & $\det[\delta^{\mu_i}_{\nu_j}]$ & generalized Kronecker delta \\
$L_k$ & \texttt{lovelockL[ra, k]} & (2.2) & Lovelock density of order $k$ ($L_0=1$, $L_1=R$) \\
$L_2$ & — & $R^2-4R_{\mu\nu}R^{\mu\nu}+R_{\mu\nu\rho\sigma}R^{\mu\nu\rho\sigma}$ & Gauss–Bonnet density \\
$L_3$ & \texttt{feL3coef} (coefficients) & (2.3) & cubic Lovelock density \\
$E^{(k)\mu}{}_\nu$ & \texttt{lovelockE}, \texttt{ECos} & (2.2) & Lovelock tensor of order $k$ ($E^{(0)}=-\delta/2$, $E^{(1)}=G$, $E^{(4)}\equiv0$) \\
$\alpha_k$ & \texttt{alf[k]} & $\alpha_0=2\Lambda$, $\alpha_1=1$ & Lovelock couplings ($\alpha_2$: length², $\alpha_3$: length\texorpdfstring{\ensuremath{{}^{4}}}{⁴}) \\
$\Lambda$ & \texttt{lamb} &  & bulk cosmological constant \\
$\kappa$ & \texttt{kapS} (sections 20–23), \texttt{k8} (section 18) & $8\pi G_8=M_8^{-6}$ & eight-dimensional gravitational coupling \\
$G_8$, $M_8$ & \texttt{rIIM8} ($M_8=10^4$ GeV) & $\kappa=8\pi G_8=M_8^{-6}$ & Newton constant and reduced Planck mass in 8 dimensions \\
$\ell_2$, $\xi$ & \texttt{rIxi}, \texttt{xiS} & $\ell_2=\sqrt{\alpha_2}$, $\xi=\alpha_3/\alpha_2^2$ & Gauss–Bonnet length; cubic-to-Gauss–Bonnet ratio ($10^{-6}$) \\
$R_{abcd}$ (test) & \texttt{feRandCurv[seed]} & $\sum_{n=1}^3(S^{(n)}_{ac}S^{(n)}_{bd}-S^{(n)}_{ad}S^{(n)}_{bc})$ & random integer algebraic curvature tensor ($S^{(n)}$ symmetric) \\
invariants & \texttt{feInvariants[RL]} & (2.3) & the eight cubic invariants of a curvature tensor \\*
\end{longtable}}
\begin{align*}
\{{}^{\lambda}_{\mu\nu}\}&=\tfrac12g^{\lambda\rho}(\partial_\mu g_{\rho\nu}+\partial_\nu g_{\rho\mu}-\partial_\rho g_{\mu\nu}),&&(2.1) \\
R^\rho{}_{\sigma\mu\nu}&=\partial_\mu\{{}^{\rho}_{\sigma\nu}\}-\partial_\nu\{{}^{\rho}_{\sigma\mu}\}+\{{}^{\rho}_{\mu\lambda}\}\{{}^{\lambda}_{\sigma\nu}\}-\{{}^{\rho}_{\nu\lambda}\}\{{}^{\lambda}_{\sigma\mu}\}, \\
R_{\sigma\nu}&=R^\rho{}_{\sigma\rho\nu},\qquad R=g^{\sigma\nu}R_{\sigma\nu}, \\
L_k&=\frac1{2^k}\delta^{\alpha_1\beta_1\cdots\alpha_k\beta_k}_{\gamma_1\delta_1\cdots\gamma_k\delta_k}R^{\gamma_1\delta_1}{}_{\alpha_1\beta_1}\cdots R^{\gamma_k\delta_k}{}_{\alpha_k\beta_k},&&(2.2) \\
E^{(k)\mu}{}_\nu&=-\frac1{2^{k+1}}\delta^{\mu\alpha_1\beta_1\cdots\alpha_k\beta_k}_{\nu\gamma_1\delta_1\cdots\gamma_k\delta_k}R^{\gamma_1\delta_1}{}_{\alpha_1\beta_1}\cdots R^{\gamma_k\delta_k}{}_{\alpha_k\beta_k}, \\
L_3&=R^3-12RR_{\mu\nu}R^{\mu\nu}+16R^\mu{}_\nu R^\nu{}_\rho R^\rho{}_\mu+24R^{\mu\rho}R^{\nu\sigma}R_{\mu\nu\rho\sigma}&&(2.3) \\
&\quad+3RR_{\mu\nu\rho\sigma}R^{\mu\nu\rho\sigma}-24R^{\mu\nu}R_\mu{}^{\rho\sigma\tau}R_{\nu\rho\sigma\tau} \\
&\quad+4R^{\mu\nu}{}_{\rho\sigma}R^{\rho\sigma}{}_{\lambda\tau}R^{\lambda\tau}{}_{\mu\nu}-8R_\mu{}^\rho{}_\nu{}^\sigma R_\rho{}^\lambda{}_\sigma{}^\tau R_\lambda{}^\mu{}_\tau{}^\nu .
\end{align*}
\Needspace{12\baselineskip}\section*{3. The fields and the terms of the action}\addcontentsline{toc}{section}{3. The fields and the terms of the action}
{\small
\begin{longtable}{>{\raggedright\arraybackslash}p{\dimexpr 0.1700\linewidth-2\tabcolsep\relax}>{\raggedright\arraybackslash}p{\dimexpr 0.2100\linewidth-2\tabcolsep\relax}>{\raggedright\arraybackslash}p{\dimexpr 0.3400\linewidth-2\tabcolsep\relax}>{\raggedright\arraybackslash}p{\dimexpr 0.2800\linewidth-2\tabcolsep\relax}}
\toprule
\textbf{symbol} & \textbf{Mathematica} & \textbf{definition} & \textbf{meaning} \\
\midrule\endfirsthead
\toprule
\textbf{symbol} & \textbf{Mathematica} & \textbf{definition} & \textbf{meaning} \\
\midrule\endhead
\endfoot
\bottomrule\endlastfoot
$S_0,\dots,S_3$ & — & (3.1) & the four Lovelock terms of gravity \\
$\psi$, $\bar\psi$ & \texttt{dirac16complex} & $\bar\psi=\psi^\dagger\Gamma_0$ & the 16-component complex field and its Dirac conjugate \\
$\Gamma_A$, $\Gamma^A$ & \texttt{Γ17[A]}, \texttt{Γ17up[A]} & (3.2) & real 16×16 Dirac matrices; $\Gamma^A=\eta^{AA}\Gamma_A$ \\
$\tau^O_i$ & \texttt{τO[i]} & $(\tau^O_i)^2=-1_8$, mutually anticommuting & real antisymmetric 8×8 matrices of the octonions \\
$e^A{}_\mu$, $e_A{}^\mu$ & \texttt{frameData[cv, e]} (section 13) & $g_{\mu\nu}=\eta_{AB}e^A{}_\mu e^B{}_\nu$ & frame (vielbein) and inverse frame \\
$\Gamma^\mu(X)$ & — & $e_A{}^\mu\Gamma^A$ & frame matrices \\
$\omega_\mu{}^A{}_B$, $\Omega_\mu$ & — & (3.3) & spin connection and its spinor form \\
$D_\mu$ & — & $\partial_\mu+\Omega_\mu$ (bulk; brane fields: (3.6)) & covariant derivative \\
$m$ & \texttt{mPs} (section 18) &  & mass of dirac16complex \\
$g_r$ & \texttt{gg[r]} (section 15) & $r=0,\dots,8$ & the nine self-couplings (dimension length\texorpdfstring{\ensuremath{{}^{6}}}{⁶}) \\
$K$, $\Theta^K$, $\eta_K$ & — & (3.4) & multi-index, Clifford product, its sign \\
$\Theta_K$ & — & $\eta_K\Theta^K$ & lowered Clifford product \\
$U(b)$, $U_0$, $K_U$, $b_0$ & \texttt{uPot}, \texttt{u0B}, \texttt{kU} & (3.5) & stabilizing potential of the auxiliary space, its minimum value, stiffness, minimum \\
$W$, $h_{\mu\nu}$ & — & $W=\lbrace y=0\rbrace$, $h_{\mu\nu}=g_{\mu\nu}(x,0)$ & world brane and its induced metric \\
$T^4$, $L_c$ & \texttt{lcT} & $x_k\simeq x_k+L_c$ & auxiliary torus and its coordinate period \\
$f_W$ & \texttt{fWb} &  & brane tension (energy per brane volume) \\
$\mathcal L_{\rm SM}$ & — & (3.6) & Standard Model Lagrangian on the brane \\
$G^a_\mu$, $W^I_\mu$, $B_\mu$ & — & (3.6) & gauge fields of $SU(3)$, $SU(2)$, $U(1)_Y$ \\
$g_s,g,g'$ & — &  & gauge couplings \\
$T^a$, $f^{abc}$, $\tau^I$, $\epsilon^{IJK}$ & — & $[T^a,T^b]=if^{abc}T^c$ & generators and structure constants \\
$\Phi$, $\tilde\Phi$ & — & $\tilde\Phi=i\sigma_2\Phi^*$ & Higgs doublet and its conjugate \\
$\mu^2$, $\lambda$ & — &  & Higgs potential parameters \\
$Y_u,Y_d,Y_e,Y_\nu$ & — &  & Yukawa matrices \\
$Q_L$, $u_R$, $d_R$, $L_L$, $e_R$, $\nu_R$, $f$ & — &  & quark and lepton fields; $f$ any of them \\
$y_{is}$, $G^f_P$, $\lambda_P$ & — & (3.7) & Yukawa, vector and Higgs portal couplings \\
$\psi_s$, $\chi_s$, $P_R$ & — & $P_R=\tfrac12(1+\gamma_5)$ & brane-right-handed parts of dirac16complex at the brane; portal spinors \\
$\phi$ & — & $\sqrt{12}\,M_{\rm Pl}\ln(b/b_0)$ & canonical radion \\
$\mathcal L_{\phi W}$ & — & (3.8) & radion–brane coupling \\
$b_{\rm QCD}$, $n_f$, $\alpha_s$ & — & $b_{\rm QCD}=11-\tfrac23n_f$ & first coefficient of the QCD beta function, active quark flavours, strong coupling \\*
\end{longtable}}
\begin{align*}
S_k&=-\frac{\alpha_k}{2\kappa}\int d^8X\sqrt{|g|}\,L_k\ \ (k=0,1,2,3),&&(3.1) \\
S_\psi&=\int d^8X\sqrt{|g|}\,\bar\psi\,i\Gamma^\mu(\partial_\mu+\Omega_\mu)\psi,\qquad S_m=-\int d^8X\sqrt{|g|}\,m\bar\psi\psi, \\
S_{\rm self}&=-\frac12\sum_{r=0}^8g_r\int d^8X\sqrt{|g|}\sum_{|K|=r}\eta_K(\bar\psi\Theta^K\psi)^2,\qquad S_U=-\int d^8X\sqrt{|g|}\,U(b), \\
S_f&=-\int_Wd^4x\sqrt{|h|}\,f_W,\qquad S_{\rm SM}=\int_Wd^4x\sqrt{|h|}\,\mathcal L_{\rm SM}, \\
\int_Wd^4x\sqrt{|h|}(\dots)&=\int d^8X\sqrt{|g|}\,\frac{\delta^4(y)}{b^4}(\dots), \\
\Gamma_0&=\begin{pmatrix}0&1_8\\1_8&0\end{pmatrix},\quad \Gamma_i=\begin{pmatrix}\tau^O_i&0\\0&-\tau^O_i\end{pmatrix},\quad \Gamma_A\Gamma_B+\Gamma_B\Gamma_A=2\eta_{AB}1_{16},&&(3.2) \\
\omega_\mu{}^A{}_B&=e^A{}_\nu\big(\partial_\mu e_B{}^\nu+\{{}^{\nu}_{\mu\lambda}\}e_B{}^\lambda\big),\qquad \Omega_\mu=\tfrac14\eta_{AC}\,\omega_\mu{}^C{}_B\,\Gamma^A\Gamma^B,&&(3.3) \\
\Theta^K&=i^{\sigma_r}\Gamma^{A_1}\cdots\Gamma^{A_r},\quad K=(A_1<\dots<A_r),\quad \sigma_r=\tfrac{r(r-1)}2\bmod2,&&(3.4) \\
\eta_K&=\eta_{A_1A_1}\cdots\eta_{A_rA_r},\qquad \Theta^{()}=1_{16}, \\
U(b)&=U_0+\frac{K_U}2\Big(\ln\frac b{b_0}\Big)^2,\qquad b^8=\det(-g_{kl}),&&(3.5) \\
\mathcal L_{\rm SM}&=-\tfrac14G^a_{\mu\nu}G^{a\mu\nu}-\tfrac14W^I_{\mu\nu}W^{I\mu\nu}-\tfrac14B_{\mu\nu}B^{\mu\nu}+(D_\mu\Phi)^\dagger D^\mu\Phi&&(3.6) \\
&\quad-\mu^2\Phi^\dagger\Phi-\lambda(\Phi^\dagger\Phi)^2+\sum_f\bar fi\gamma^\mu D_\mu f \\
&\quad-\big[\bar Q_LY_u\tilde\Phi u_R+\bar Q_LY_d\Phi d_R+\bar L_LY_e\Phi e_R+\bar L_LY_\nu\tilde\Phi\nu_R+\text{h.c.}\big], \\
D_\mu&=\nabla_\mu-ig_sG^a_\mu T^a-igW^I_\mu\tfrac{\tau^I}2-ig^\prime YB_\mu, \\
G^a_{\mu\nu}&=\partial_\mu G^a_\nu-\partial_\nu G^a_\mu+g_sf^{abc}G^b_\mu G^c_\nu,\qquad W^I_{\mu\nu}=\partial_\mu W^I_\nu-\partial_\nu W^I_\mu+g\epsilon^{IJK}W^J_\mu W^K_\nu, \\
B_{\mu\nu}&=\partial_\mu B_\nu-\partial_\nu B_\mu, \\
S_P&=\int_Wd^4x\sqrt{|h|}\Big[-\sum_{i,s}y_{is}\bar L_i\tilde\Phi P_R\psi_s+\text{h.c.}&&(3.7) \\
&\qquad-G^f_P(\bar\psi\Gamma^\mu\psi)|_W(\bar f\gamma_\mu f)-\lambda_P\Phi^\dagger\Phi(\bar\psi\psi)|_W\Big], \\
\mathcal L_{\phi W}&=\frac{\phi}{\sqrt3M_{\rm Pl}}T^\mu{}_\mu\big|_{\rm brane},\qquad T^\mu{}_\mu\supset m_f\bar ff,\qquad T^\mu{}_\mu\supset-\frac{b_{\rm QCD}\alpha_s}{8\pi}G^a_{\mu\nu}G^{a\mu\nu}.&&(3.8)
\end{align*}
\Needspace{12\baselineskip}\section*{4. Energy–momentum tensors and the homogeneous metric}\addcontentsline{toc}{section}{4. Energy–momentum tensors and the homogeneous metric}
{\small
\begin{longtable}{>{\raggedright\arraybackslash}p{\dimexpr 0.1700\linewidth-2\tabcolsep\relax}>{\raggedright\arraybackslash}p{\dimexpr 0.2100\linewidth-2\tabcolsep\relax}>{\raggedright\arraybackslash}p{\dimexpr 0.3400\linewidth-2\tabcolsep\relax}>{\raggedright\arraybackslash}p{\dimexpr 0.2800\linewidth-2\tabcolsep\relax}}
\toprule
\textbf{symbol} & \textbf{Mathematica} & \textbf{definition} & \textbf{meaning} \\
\midrule\endfirsthead
\toprule
\textbf{symbol} & \textbf{Mathematica} & \textbf{definition} & \textbf{meaning} \\
\midrule\endhead
\endfoot
\bottomrule\endlastfoot
$T_{\mu\nu}$ & — & $\frac2{\sqrt{|g|}}\frac{\delta S}{\delta g^{\mu\nu}}$ & energy–momentum tensor of a sector ($T^0{}_0=\varepsilon$) \\
$T_\psi$, $T_U$, $T_f$, $T_{\rm SM}$, $T_P$ & — &  & energy–momentum tensors of the sectors \\
$a$, $b$, $n$ & \texttt{aS}, \texttt{bS}; \texttt{a3}, \texttt{b4} & (4.1) & scale factors of ordinary and auxiliary space; lapse ($n=1$) \\
$H_a$, $H_b$ & \texttt{Ha}, \texttt{Hb} & $a'/a$, $b'/b$ & Hubble rates \\
$H_a'$, $H_b'$ & \texttt{dHa}, \texttt{dHb} & $dH_a/dx_0$, $dH_b/dx_0$ & their derivatives \\
$\varepsilon$, $p_a$, $p_b$ & \texttt{epsM}, \texttt{paM}, \texttt{pbM} & $T^\mu{}_\nu=\mathrm{diag}(\varepsilon,-p_a\times3,-p_b\times4)$ & energy density, pressures along ordinary and auxiliary space \\
$W(a,b)$ & \texttt{wW[a, b]} & $a^3b^4\varepsilon$ & comoving energy at fixed occupation numbers \\
$k_a$, $k_b$ & \texttt{kaP}, \texttt{kbP} &  & comoving momentum of a particle along ordinary/auxiliary space \\
$\omega$ & \texttt{root} in \texttt{rIgas} & (4.3) & energy of a massless particle (times $a^3b^4/V$) \\
$V_\perp$ & \texttt{vPerpS}, \texttt{rIIVp} & $b^4L_c^4$ & proper volume of the auxiliary torus \\
$\rho_4$, $p_4$, $w$ & \texttt{rho4}, \texttt{wW0} & $p_4=w\rho_4$ & energy density and pressure of a brane fluid \\
$\theta$, $\langle\cdot\rangle$ & \texttt{rIth}, \texttt{rIwt} & (4.3) & momentum angle and the average over the momentum sphere \\
$C_g$ & \texttt{rICg} & $3\times10^{-6}$ (units $\ell_2=1$) & normalization of the dirac16complex gas \\
$\mathrm{EL}_n,\mathrm{EL}_a,\mathrm{EL}_b$ & \texttt{feELn}, \texttt{feELa}, \texttt{feELb} & Euler–Lagrange expressions & minisuperspace variations \\
$e_{00}$, $e_{11}$, $e_{44}$ & \texttt{eTot[1]}, \texttt{eTot[2]}, \texttt{eTot[5]}; \texttt{feTot[c]} ($\alpha_1=1$) & $\sum_k\alpha_kE^{(k)c}{}_c-\Lambda$ & left-hand sides of the 00, 11, 44 equations \\*
\end{longtable}}
\begin{align*}
ds^2&=n^2dx_0^2-a^2(dx_1^2+dx_2^2+dx_3^2)-b^2(dx_4^2+\dots+dx_7^2),&&(4.1) \\
\varepsilon&=\frac W{a^3b^4},\qquad p_a=-\frac{\partial_aW}{3a^2b^4},\qquad p_b=-\frac{\partial_bW}{4a^3b^3},&&(4.2) \\
0&=\varepsilon'+3H_a(\varepsilon+p_a)+4H_b(\varepsilon+p_b), \\
\kappa\varepsilon&=\frac{C_g}{a^3b^4}\langle\omega\rangle,\qquad \kappa p_a=\frac{C_g}{a^3b^4}\Big\langle\frac{\cos^2\theta}{3a^2\omega}\Big\rangle,\qquad \kappa p_b=\frac{C_g}{a^3b^4}\Big\langle\frac{\sin^2\theta}{4b^2\omega}\Big\rangle,&&(4.3) \\
\omega&=\sqrt{\frac{\cos^2\theta}{a^2}+\frac{\sin^2\theta}{b^2}},\qquad \langle f\rangle=\frac{\int_0^{\pi/2}f\cos^2\theta\sin^3\theta\,d\theta}{\int_0^{\pi/2}\cos^2\theta\sin^3\theta\,d\theta}.
\end{align*}
The field equations of the homogeneous metric (00, 11, 44 components) and the 00 equation written out:
\begin{align*}
\sum_k\alpha_kE^{(k)0}{}_0&=\kappa\varepsilon,\qquad \sum_k\alpha_kE^{(k)1}{}_1=-\kappa p_a,\qquad \sum_k\alpha_kE^{(k)4}{}_4=-\kappa p_b,&&(4.4) \\
\kappa\varepsilon&=3(H_a^2+4H_aH_b+2H_b^2)-12\alpha_2H_b(4H_a^3+18H_a^2H_b+12H_aH_b^2+H_b^3) \\
&\quad+360\alpha_3H_a^2H_b^3(4H_a+3H_b)-\Lambda .
\end{align*}
\Needspace{12\baselineskip}\section*{5. The radion and the four-dimensional constants}\addcontentsline{toc}{section}{5. The radion and the four-dimensional constants}
{\small
\begin{longtable}{>{\raggedright\arraybackslash}p{\dimexpr 0.1700\linewidth-2\tabcolsep\relax}>{\raggedright\arraybackslash}p{\dimexpr 0.2100\linewidth-2\tabcolsep\relax}>{\raggedright\arraybackslash}p{\dimexpr 0.3400\linewidth-2\tabcolsep\relax}>{\raggedright\arraybackslash}p{\dimexpr 0.2800\linewidth-2\tabcolsep\relax}}
\toprule
\textbf{symbol} & \textbf{Mathematica} & \textbf{definition} & \textbf{meaning} \\
\midrule\endfirsthead
\toprule
\textbf{symbol} & \textbf{Mathematica} & \textbf{definition} & \textbf{meaning} \\
\midrule\endhead
\endfoot
\bottomrule\endlastfoot
$\beta$ & \texttt{betR}; \texttt{lb - rIIlnb0} & $\ln(b/b_0)$ & radion (logarithmic size of the auxiliary space), $H_b=\beta'$ \\
$m_r$ & \texttt{rIIm} ($10^{-3}/\ell_2$), \texttt{rIIImr} & $m_r^2=\kappa K_U/12$ & radion mass \\
$\beta_{\rm eq}$ & \texttt{bqV} & (5.2) & equilibrium of the radion \\
$G_4$, $\Lambda_4$ & — & (5.1) & four-dimensional Newton constant and cosmological constant \\
$M_{\rm Pl}$ & \texttt{rIIMPl} & $2.435\times10^{18}$ GeV, $M_{\rm Pl}^2=M_8^6V_\perp$ & reduced four-dimensional Planck mass \\
$a_E$, $H_E$ & \texttt{rIIHE} & $a_E=ab^2$, $H_E=H_a+2H_b$ & Einstein-frame scale factor and expansion rate \\
$g_E$ & — & $h=(b_0/b)^4g_E$ & Einstein-frame metric of the brane \\
$\ell_\perp$ & \texttt{rIIVp\textasciicircum{}(1/4)} & $V_\perp^{1/4}$ & size of the auxiliary torus ($3.08\times10^{-13}$ m) \\
$m_f$, $N_c$ & — &  & fermion mass, number of colours \\
$\Gamma$, $T_{\rm RH}$ & \texttt{rIIgam}, \texttt{rIITRH} & (5.3) & radion decay rate, reheating temperature \\
$N_{\rm KK}$ & \texttt{rIINKK} & $V_\perp\frac{\pi^2}2\frac{(m_r/2)^4}{(2\pi)^4}$ & Kaluza–Klein momenta below $m_r/2$ \\
$T_*$ & — & $(M_8^6/M_{\rm Pl})^{1/5}$ & normalcy temperature (estimate) \\*
\end{longtable}}
\begin{align*}
3H_a^2&=8\pi G_4\rho_4+\Lambda_4,\qquad H_a'=-\frac\kappa2(\varepsilon+p_a),&&(5.1) \\
8\pi G_4&=\frac\kappa{V_\perp},\qquad \Lambda_4=\Lambda+\kappa U_0+\frac{\kappa f_W}{V_\perp}, \\
\beta''+3H_a\beta'+4\beta'^2&=\kappa\Big(\frac\varepsilon6-\frac{p_a}2+\frac{p_b}3\Big)+\frac\Lambda3\ \ (\alpha_2=\alpha_3=0),\qquad 3H_E^2=6\beta'^2+\kappa\varepsilon, \\
\frac\kappa4bU'(b)&=\Lambda+\kappa U+\frac{2\kappa f_W}{V_\perp}+\frac\kappa2(\varepsilon-3p_a)+\kappa p_b-12\alpha_2H_a^2(H_a'+H_a^2)\quad(\text{static }b),&&(5.2) \\
\beta_{\rm eq}&=\frac{[\dots]}{3m_r^2}, \\
\Gamma(\phi\to f\bar f)&=\frac{N_cm_f^2m_r}{24\pi M_{\rm Pl}^2}\Big(1-\frac{4m_f^2}{m_r^2}\Big)^{3/2},\qquad \Gamma(\phi\to gg)=\frac{b_{\rm QCD}^2\alpha_s^2m_r^3}{96\pi^3M_{\rm Pl}^2},&&(5.3) \\
T_{\rm RH}&=\Big[\frac{\Gamma M_{\rm Pl}}{(\pi^2g_*/90)^{1/2}}\Big]^{1/2}.
\end{align*}
In (5.2) $[\dots]$ is the right-hand side of the equation for $\tfrac\kappa4bU'(b)$ with $U$ replaced by $U_0$; with $U=U_0+\tfrac{K_U}2\beta^2$ one has $bU'=K_U\beta$ and $\tfrac\kappa4K_U=3m_r^2$.

\Needspace{12\baselineskip}\section*{6. Regime I: the primordial Lovelock regime}\addcontentsline{toc}{section}{6. Regime I: the primordial Lovelock regime}
{\small
\begin{longtable}{>{\raggedright\arraybackslash}p{\dimexpr 0.1700\linewidth-2\tabcolsep\relax}>{\raggedright\arraybackslash}p{\dimexpr 0.2100\linewidth-2\tabcolsep\relax}>{\raggedright\arraybackslash}p{\dimexpr 0.3400\linewidth-2\tabcolsep\relax}>{\raggedright\arraybackslash}p{\dimexpr 0.2800\linewidth-2\tabcolsep\relax}}
\toprule
\textbf{symbol} & \textbf{Mathematica} & \textbf{definition} & \textbf{meaning} \\
\midrule\endfirsthead
\toprule
\textbf{symbol} & \textbf{Mathematica} & \textbf{definition} & \textbf{meaning} \\
\midrule\endhead
\endfoot
\bottomrule\endlastfoot
$\tau$, $x_s$ & \texttt{tS}; \texttt{rIxs} & $\tau=x_0-x_s$ & time since the singularity; its position \\
$u$, $X$ & \texttt{uS}, \texttt{xS} & $u=H_b/H_a$, $X=H_a^2$ & variables of the vacuum curve \\
$P_E,P_G,P_C$ & \texttt{rIpE}, \texttt{rIpG}, \texttt{rIpC} & (6.1) & Einstein, Gauss–Bonnet, cubic parts of the vacuum constraint \\
$X(u)$ & \texttt{rIX[u, xi]} & (6.1) & the vacuum curve \\
$F$ & \texttt{rIF[u, xi]} & $-H_a'$ on the curve & rate along the curve \\
$p_{\rm eff}$ & \texttt{rIpeff[u, xi]}, \texttt{"peff"} & $-H_a^2/H_a'$ & effective Kasner exponent ($=p$ for $a\propto\tau^p$) \\
$u_1,p_1,q_1$ & \texttt{rIu1}, \texttt{rIp1}, \texttt{rIq1} & (6.2) & Einstein–Kasner constants \\
$u_2,p_2,q_2$ & \texttt{rIu2}, \texttt{rIp2}, \texttt{rIq2} & (6.2) & Gauss–Bonnet–Kasner constants \\
$h_0,h_1,h_2,A_0,A_1$ & \texttt{rIcub}; \texttt{ch1}, \texttt{ch2}, \texttt{cA0}, \texttt{cA1} & (6.3) & coefficients of the cubic-regime series \\
$d_1$, $e_1$ & \texttt{rId1}, \texttt{rIe1} & (6.4) & Einstein-regime coefficients \\
$f_1$, $f_2$ & \texttt{rIf1}, \texttt{rIf2} & $p_{\rm eff}=p_2+f_1/H_a^2+f_2\xi H_a^2$ & Gauss–Bonnet-regime coefficients \\
$s_1$, $s_2$ & \texttt{rIs1v}, \texttt{rIs2v} & (6.4), (6.5) & time shifts of the Einstein and Gauss–Bonnet formulas \\
$H_{a,\rm mid}$ & \texttt{rIHmid} & $(f_1/(f_2\xi))^{1/4}$ & middle of the Gauss–Bonnet regime \\
$\tau(u_2)$, $H_a(u_2)$, $J$ & \texttt{rItauU2}, \texttt{rIHu2}, \texttt{rIJ} & (6.4) & quadrature data \\
$\tau_{32}$, $\tau_{21}$ & — & $|q_2|\sqrt{30\xi}\,\ell_2$, $\ell_2$ & crossover times \\
$h_*$ & \texttt{rIhStar} & $-\sqrt{\alpha_2/(54\alpha_3)}$ & auxiliary contraction rate with matter near $\tau=0$ \\
$w_a,w_b$ & \texttt{wA} & $p_a=w_a\varepsilon$, $p_b=w_b\varepsilon$ & equations of state \\
$P$, $Q$, $\mathcal C$ & \texttt{pP}, \texttt{qQ}, \texttt{cC} & $a\propto\tau^P$, $b\propto\tau^Q$, $\mathcal C=\kappa\varepsilon\tau^2$ & power laws of Einstein gravity with a fluid \\
$S_1,S_2,S_3$ & \texttt{"S"} in \texttt{rIstate} & (6.6) & sizes of the curvature terms of the 00 equation \\
$f_{\rm gas}$ & \texttt{"frac"} & (6.6) & gas fraction \\
$K$, $F(K)$ & \texttt{rIbog[k, p, q]}, \texttt{rIFk} & $|\beta_K|^2$ & scaling variable of a mode; Bogoliubov coefficient (pair creation) \\
$\sigma_1,\sigma_2$ & \texttt{rIsg1}, \texttt{rIsg2} & Pauli matrices & mode Hamiltonian $K\tau^{-p_1}\sigma_1+K\tau^{-q_1}\sigma_2$ \\
$\sigma$ & \texttt{sV} & $-\ln H_a$ & integration variable \\*
\end{longtable}}
\begin{align*}
P_E&=3(2u^2+4u+1),\quad P_G=-12u(u^3+12u^2+18u+4),\quad P_C=360u^3(3u+4),&&(6.1) \\
X(u)&=\frac{P_G+\sqrt{P_G^2-4\xi P_CP_E}}{2\xi|P_C|}, \\
u_1&=-1+\tfrac1{\sqrt2},\quad p_1=\frac1{3+4u_1}=\frac{1+2\sqrt2}7,\quad q_1=u_1p_1,&&(6.2) \\
0&=u_2^3+12u_2^2+18u_2+4,\ \ u_2\in(-\tfrac13,0),\quad p_2=\frac3{3+4u_2},\quad q_2=u_2p_2, \\
H_a&=\frac1\tau+A_0+A_1\tau,\quad H_b=h_0+h_1\tau+h_2\tau^2,\quad h_0=-\sqrt{\frac{\alpha_2}{30\alpha_3}},\quad h_1=\frac{\alpha_2}{16\alpha_3}-\frac1{32\alpha_2},&&(6.3) \\
A_0&=\frac{\sqrt{30}(58\alpha_2^2-45\alpha_3)}{1920\alpha_2^{3/2}\alpha_3^{1/2}},\qquad A_1=-\frac{3647\alpha_2}{92160\alpha_3}-\frac{223}{6144\alpha_2}+\frac{645\alpha_3}{8192\alpha_2^3}, \\
h_2&=\frac{\sqrt{30}(-1212\alpha_2^4+1180\alpha_2^2\alpha_3+225\alpha_3^2)}{307200\alpha_2^{5/2}\alpha_3^{3/2}}, \\
u&=u_1+d_1H_a^2,\ d_1=\frac{\alpha_2(16-11\sqrt2)}4,\quad p_{\rm eff}=p_1+e_1H_a^2,\ e_1=\frac{\alpha_2(1+9\sqrt2)}7,&&(6.4) \\
\tau&=\frac{p_1}{H_a}-e_1H_a+s_1,\qquad s_1=\tau(u_2)-\frac{p_1}{H_a(u_2)}+J, \\
J&=\int_0^{H_a(u_2)}\Big(\frac1F-\frac{p_1}{H^2}\Big)dH,\qquad \tau(H_a)=\int_{H_a}^\infty\frac{dH}F, \\
\tau&=\frac{p_2}{H_a}+\frac{f_1}{3H_a^3}-f_2\xi H_a+s_2\quad\text{(Gauss–Bonnet regime)},&&(6.5) \\
S_1&=3(H_a^2+4H_a|H_b|+2H_b^2),\qquad S_2=12\alpha_2|H_b|(4H_a^3+18H_a^2|H_b|+12H_aH_b^2+|H_b|^3),&&(6.6) \\
S_3&=360\alpha_3H_a^2|H_b|^3(4H_a+3|H_b|),\qquad f_{\rm gas}=\frac{\kappa\varepsilon}{S_1+S_2+S_3+\kappa\varepsilon}.
\end{align*}
\Needspace{12\baselineskip}\section*{7. Regime II: the modern regime}\addcontentsline{toc}{section}{7. Regime II: the modern regime}
{\small
\begin{longtable}{>{\raggedright\arraybackslash}p{\dimexpr 0.1700\linewidth-2\tabcolsep\relax}>{\raggedright\arraybackslash}p{\dimexpr 0.2100\linewidth-2\tabcolsep\relax}>{\raggedright\arraybackslash}p{\dimexpr 0.3400\linewidth-2\tabcolsep\relax}>{\raggedright\arraybackslash}p{\dimexpr 0.2800\linewidth-2\tabcolsep\relax}}
\toprule
\textbf{symbol} & \textbf{Mathematica} & \textbf{definition} & \textbf{meaning} \\
\midrule\endfirsthead
\toprule
\textbf{symbol} & \textbf{Mathematica} & \textbf{definition} & \textbf{meaning} \\
\midrule\endhead
\endfoot
\bottomrule\endlastfoot
$\tau_E$, $A$, $\vartheta$ & \texttt{rIItEF}, \texttt{rIIaF}, \texttt{rIIphF} & (7.1) & Einstein-frame time origin, amplitude and phase of the radion oscillation \\
$\ln b_0$ & \texttt{rIIlnb0} &  & minimum of $U$ in the capture run \\
$T$, $T_1$ & — & $T_1=10$ GeV & photon temperature; start of the $t(T)$ integration \\
$g_*$, $g_{*s}$ & \texttt{rIIgstars[T]} & (7.2) & relativistic degrees of freedom for energy and entropy \\
$g_i$, $m_i$, $x$ & \texttt{rIIsp} & $x=m_i/T$ & degeneracy and mass of species $i$ \\
$K_1$, $K_2$ & \texttt{BesselK[1, …]}, \texttt{BesselK[2, …]} &  & modified Bessel functions of the second kind \\
$s$ & — & $(2\pi^2/45)g_{*s}T^3$ & entropy density \\
$S(T)$ & \texttt{rIIS[T]} & $\tfrac12[1+\tanh((T-0.160)/0.030)]$ (GeV) & QCD switch (quarks and gluons \texorpdfstring{\ensuremath{\leftrightarrow}}{↔} hadrons) \\
$T_\nu$, $N_{\rm eff}$, $g_{\gamma e}$ & \texttt{rIIgge} & $T_\nu=T\min(1,(g_{\gamma e}/\tfrac{11}2)^{1/3})$, $N_{\rm eff}=3.044$ & neutrino temperature, effective number, $g_{*s}$ of $\gamma$ and $e^\pm$ \\
$t(T)$ & \texttt{rIItg} & (7.3) & time–temperature relation of the radiation era \\
$H_0$, $\Omega_m$, $\Omega_r$, $\Omega_\Lambda$, $\Omega_\gamma$ & \texttt{rIIH0s}, \texttt{rIIOm}, \texttt{rIIOr}, \texttt{rIIOL}, \texttt{rIIOg} & (7.4) & Planck 2018 parameters and the radiation densities \\
$T_{\rm CMB}$ & \texttt{rIITcmb} & 2.7255 K & temperature of the cosmic microwave background \\
$t_{rm}$, $t_{m\Lambda}$, $t_m$ & \texttt{rIItrm}, \texttt{rIItmL}, \texttt{rIItm} & (7.4) & exact two-component solutions and the matter solution \\
$z_{\rm eq}$, $t_0$ & \texttt{rIIage} & $\Omega_m/\Omega_r-1$; $\int_0^1\frac{da}{aH_a}$ & redshift of matter–radiation equality, age \\*
\end{longtable}}
\begin{align*}
\beta&\approx\frac{A\sin(m_r\tau+\vartheta)}{m_r(\tau-\tau_E)},\qquad H_E\approx\frac2{3(\tau-\tau_E)},&&(7.1) \\
\rho_i&=\frac{g_iT^4}{2\pi^2}\sum_{n\ge1}(\pm1)^{n+1}\frac{3(nx)^2K_2(nx)+(nx)^3K_1(nx)}{n^4},&&(7.2) \\
p_i&=\frac{g_iT^4}{2\pi^2}\sum_{n\ge1}(\pm1)^{n+1}\frac{(nx)^2K_2(nx)}{n^4},\qquad g_*=\frac{30\rho}{\pi^2T^4},\qquad g_{*s}=\frac{45s}{2\pi^2T^3}, \\
H_a&=\Big(\frac{\pi^2g_*}{90}\Big)^{1/2}\frac{T^2}{M_{\rm Pl}},\qquad \ln a=-\ln T-\tfrac13\ln g_{*s}+\text{const},&&(7.3) \\
t(T)&=t(T_1)+\int\frac{d\ln a}{H_a},\qquad t(T_1)=\frac1{2H_a(T_1)}, \\
H_a&=H_0\sqrt{\Omega_ra^{-4}+\Omega_ma^{-3}+\Omega_\Lambda},\qquad \Omega_\gamma=\frac{(\pi^2/15)T_{\rm CMB}^4}{3H_0^2M_{\rm Pl}^2},&&(7.4) \\
\Omega_r&=\Omega_\gamma\Big[1+\frac78\Big(\frac4{11}\Big)^{4/3}N_{\rm eff}\Big],\qquad t_m=\frac{2a^{3/2}}{3H_0\Omega_m^{1/2}}, \\
t_{rm}&=\frac2{3H_0\Omega_m^2}\Big[(\Omega_ma-2\Omega_r)\sqrt{\Omega_r+\Omega_ma}+2\Omega_r^{3/2}\Big],\qquad t_{m\Lambda}=\frac2{3H_0\Omega_\Lambda^{1/2}}\mathrm{arcsinh}\Big(\sqrt{\frac{\Omega_\Lambda}{\Omega_m}}a^{3/2}\Big).
\end{align*}
In (7.2) the upper sign is for bosons, the lower (alternating) for fermions; $s_i=(\rho_i+p_i)/T$.

\Needspace{12\baselineskip}\section*{8. Regime III: the far future}\addcontentsline{toc}{section}{8. Regime III: the far future}
{\small
\begin{longtable}{>{\raggedright\arraybackslash}p{\dimexpr 0.1700\linewidth-2\tabcolsep\relax}>{\raggedright\arraybackslash}p{\dimexpr 0.2100\linewidth-2\tabcolsep\relax}>{\raggedright\arraybackslash}p{\dimexpr 0.3400\linewidth-2\tabcolsep\relax}>{\raggedright\arraybackslash}p{\dimexpr 0.2800\linewidth-2\tabcolsep\relax}}
\toprule
\textbf{symbol} & \textbf{Mathematica} & \textbf{definition} & \textbf{meaning} \\
\midrule\endfirsthead
\toprule
\textbf{symbol} & \textbf{Mathematica} & \textbf{definition} & \textbf{meaning} \\
\midrule\endhead
\endfoot
\bottomrule\endlastfoot
$H_\infty$ & \texttt{rIIIH} & $H_0\Omega_\Lambda^{1/2}$ & final de Sitter expansion rate \\
$T_{\rm GH}$ & — & $H_\infty/(2\pi)$ & Gibbons–Hawking temperature \\
$M$ & — &  & mass of a created particle \\
$\gamma$ & \texttt{rIIIgam} & (8.1) & friction coefficient of the radion in de Sitter \\
$m_{\rm eff}$ & \texttt{rIIImeff} & (8.1) & effective radion mass in de Sitter \\
$H$ & \texttt{hS}, \texttt{h8} &  & de Sitter rate of the linearization; of eight-dimensional de Sitter \\
$\Lambda$, $m_r$ (test) & \texttt{rIIIlam}, \texttt{rIIImr} & $3\times10^{-4}$, $0.1$ (units $\ell_2=1$) & parameters of the numerical de Sitter test \\*
\end{longtable}}
\begin{align*}
a(t)&=\Big(\frac{\Omega_m}{\Omega_\Lambda}\Big)^{1/3}\sinh^{2/3}\Big(\frac32H_\infty t\Big),\qquad H_a=H_\infty\coth\Big(\frac32H_\infty t\Big),\qquad |\beta_M|^2=\frac1{e^{2\pi M/H_\infty}+1}, \\
0&=\beta''+\gamma\beta'+m_{\rm eff}^2(\beta-\beta_{\rm eq}),&&(8.1) \\
\frac\gamma{3H}&=\frac{1-16\alpha_2H^2+48\alpha_2^2H^4-48\alpha_3H^4+512\alpha_2^3H^6+192\alpha_2\alpha_3H^6-1024\alpha_2^4H^8}{(1-4\alpha_2H^2+32\alpha_2^2H^4)^2}, \\
\frac{m_{\rm eff}^2}{m_r^2}&=\frac1{1-4\alpha_2H^2+32\alpha_2^2H^4},\qquad 3m_r^2\beta_{\rm eq}=3H^2-12\alpha_2H^4, \\
\Lambda&=21H^2-420\alpha_2H^4+2520\alpha_3H^6\quad(a=b=e^{Hx_0}).
\end{align*}
\Needspace{12\baselineskip}\section*{9. Symbols that look alike}\addcontentsline{toc}{section}{9. Symbols that look alike}
{\small
\begin{longtable}{>{\raggedright\arraybackslash}p{\dimexpr 0.4000\linewidth-2\tabcolsep\relax}>{\raggedright\arraybackslash}p{\dimexpr 0.6000\linewidth-2\tabcolsep\relax}}
\toprule
\textbf{symbol} & \textbf{not to be confused with} \\
\midrule\endfirsthead
\toprule
\textbf{symbol} & \textbf{not to be confused with} \\
\midrule\endhead
\endfoot
\bottomrule\endlastfoot
$\beta=\ln(b/b_0)$ (radion) & $|\beta_K|^2=F(K)$, $|\beta_M|^2$ (Bogoliubov coefficients); $\alpha_k,\beta_k$ (dummy indices of the generalized delta) \\
$\phi$ (canonical radion) & $\vartheta$ (phase of the radion oscillation); $\Phi$ (Higgs doublet); $\theta$ (momentum angle) \\
$b$ (auxiliary scale factor), $b_0$ (minimum of $U$) & $b_{\rm QCD}$ (QCD beta-function coefficient); the $b$ quark \\
$\Gamma_A$ (16×16 Dirac matrices), $\Gamma^\mu(X)$ (frame matrices) & $\Gamma$ (decay rate); $\{{}^{\lambda}_{\mu\nu}\}$ (Christoffel symbols) \\
$\gamma$ (radion friction) & $\gamma^\mu$, $\gamma_5$ (4×4 Dirac matrices of the brane) \\
$S_0,\dots,S_3,S_\psi,\dots$ (terms of the action) & $S_1,S_2,S_3$ (sizes of the curvature terms, (6.6)); $S(T)$ (QCD switch); $s$ (entropy density); $s_1,s_2$ (time shifts) \\
$W$ (world brane) & $W(a,b)$ (comoving energy); $W^I_\mu$ (electroweak gauge fields) \\
$C_g$ (gas normalization) & $\mathcal C=\kappa\varepsilon\tau^2$ (power-law constant); $N_c$ (colours) \\
$F$ (rate $-H_a'$ along the vacuum curve) & $F(K)$ (Bogoliubov coefficient) \\
$X$ (coordinates $x_0,\dots,x_7$) & $X=H_a^2$ (variable of the vacuum curve) \\
$\eta$ (frame metric: the squares of the octonion basis) & $\eta_K$ (sign of a multi-index); $\{7,1,0\}$ (xAct signature list: seven space-like, one time-like dimension) \\
$\kappa$ ($8\pi G_8$) & $K$ (scaling variable of a mode, or a multi-index); $K_U$ (stiffness of $U$); $K_1,K_2$ (Bessel functions) \\
$p_1,p_2,q_1,q_2$ (Kasner exponents) & $p_a,p_b$ (pressures); $p_{\rm eff}$; $P$, $Q$ (fluid power laws); $P_E,P_G,P_C$ (polynomials) \\
$h_0$ (vacuum cubic-regime rate) & $h_*$ (with matter); $h_{\mu\nu}$ (induced brane metric) \\
$\Lambda$ (bulk) & $\Lambda_4$ (four-dimensional); $\Omega_\Lambda$ \\
$\tau$ (time since the singularity) & $\tau^O_i$ (octonion matrices), $\tau^I$ (Pauli matrices of $SU(2)$), $\tau$ lepton; $\tau_E$, $\tau_{32}$, $\tau_{21}$ \\*
\end{longtable}}
\end{document}
